% Part: computability
% Chapter: computability-theory
% Section: coding-computations

\documentclass[../../../include/open-logic-section]{subfiles}

\begin{document}

\olfileid{cmp}{thy}{cod}
\olsection{గణనలను సంకేతీకరించడం}

ప్రతి గణనా నమూనాలోనూ కింది పనులు చేయవచ్చు:
\begin{enumerate}
\item గణనీయ ప్రమేయాల \emph{నిర్వచనాలను} క్రమబద్ధంగా వర్ణించడం.
  ఉదాహరణకు, ట్యూరింగ్ యంత్ర నిర్దేశాలు, పునరావృత్త నిర్వచనాలు, లేదా
  ఒక ప్రోగ్రామింగ్ భాషలోని ప్రోగ్రాములు ఈ నిర్వచనాలను ఇస్తాయని
  భావించవచ్చు.
\item ఏదో నిర్వచనం ఇచ్చే ప్రమేయాన్ని, ఇచ్చిన ఇన్‌పుట్‌పై గణించిన
  పూర్తి లేఖనాన్ని వర్ణించడం. ఉదాహరణకు, గణనలోని ప్రతి దశకు సంబంధించిన
  స్థితివిన్యాసాల అనుక్రమంతో (యంత్ర స్థితి, టేపు విషయం) ట్యూరింగ్
  యంత్ర గణనను వర్ణించవచ్చు.
\item ఒక గణనకు సంబంధించినదిగా ప్రతిపాదించిన లేఖనం, ఇచ్చిన
  నిర్వచనం గల గణనీయ ప్రమేయాన్ని ఇచ్చిన ఇన్‌పుట్‌పై గణించిన విధానపు
  లేఖనమేనా అని పరీక్షించడం (ఆ ఇన్‌పుట్ వద్ద ప్రమేయం నిర్వచితమై,
  అంటే గణన ఆగే సందర్భంలో).
\item ఒక గణన పూర్తి లేఖనపు అటువంటి వివరణ నుంచి, ఇచ్చిన ఇన్‌పుట్‌కు
  ప్రమేయ విలువను వెలికితీయడం. ఉదాహరణకు, ఆగే ట్యూరింగ్ యంత్ర గణనలోని
  చివరి దశలో టేపుపై ఉన్న విషయమే విలువ.
\end{enumerate}

సంకేతీకరణను వాడి, ప్రతి గణనీయ ప్రమేయ వివరణకూ ఒక సంఖ్యాత్మక
\emph{సూచిక}ను కేటాయించవచ్చు; ఆ సూచిక నుంచి ఆదేశాలను గణనీయ పద్ధతిలో
తిరిగి పొందగలిగేలా ఈ కేటాయింపు ఉంటుంది. అదే విధంగా ఒక గణన పూర్తి
లేఖనాన్ని కూడా ఒకే సంఖ్యతో సంకేతీకరించవచ్చు. అప్పుడు ``సూచిక~$e$ గల
ప్రమేయాన్ని ఇన్‌పుట్~$x$పై గణించిన లేఖనాన్ని $s$ సంకేతీకరిస్తుంది'' అనే
అంకగణిత సంబంధమూ, ``సంకేతం~$s$ గల గణనా అనుక్రమపు అవుట్‌పుట్'' అనే
ప్రమేయమూ గణనీయాలు; నిజానికి అవి ఆదిమ పునరావృత్తాలు.

ఈ మౌలిక వాస్తవం చాలా శక్తివంతమైనది. ఎంచుకున్న గణనా నమూనాపై ఆధారపడకుండా,
గణనీయత గురించి అనేక ఆశ్చర్యకరమైన, ముఖ్యమైన ఫలితాలను నిరూపించడానికి అది
మనకు వీలు కల్పిస్తుంది.

\end{document}
